% Patch record for C-57 -- second-pass simplification (proposer-vs-critic agent debate).
% Applied 2026-10-28 (Sebas) to introduction.tex, methodology.tex, results.tex, conclusions.tex, on top of the
% uncommitted R3b-01 edits. Source of every change: intake/pending/second_pass_simplification_debate.md
% ("## Surviving changes" -- critic's final text used for every ACCEPT-MODIFIED item).
% Surviving: P-01 (non-ablation sites), P-02, P-04, P-05, P-06, P-07 (prose), P-08, P-09, P-10 (filler only),
%   P-11, P-12, P-13 (prose), P-14 (prose), P-16, P-17, P-18, P-19, P-20, P-21, P-22, P-23 (no caption
%   fragment; simulated GRAPH_IMU2 sentence kept; physical merge cites both panels), P-26, P-27, P-28, P-29, P-30.
% Small deviations from the debate text, both to avoid information/citation loss:
%   - P-04: kept "r in Eq. L4 runs over all auxiliary units" (the only place the index r is explained).
%   - P-23: merged physical sentence cites fig:GRAPH_IMU_real and fig:GRAPH_IMU2_real (not the parent
%     fig:graph_imu_real_pair) so both panels keep an in-text citation; Z0V2 sentence drops its table
%     pointer because the P-28 where-clause directly below already cites table:Locom_MLP.
% Rejected: P-24. Deferred to owners: P-03 (Sam R2b-01 / Deivi R3b-02), P-15 (Deivi), P-25 (Diego), ablation
%   reviewer-request clauses and "(R3-01)"/"(R3-02)" (Sam), P-07 subcaption + P-14 captions (Deivi), P-10
%   "every retained trial" consolidation (Diego).
% Inconsistencies S-01..S-10 are documented in the debate file and were NOT changed by this patch.
% Words: introduction 1532->1465, methodology 6408->5862, results 12589->11778, conclusions 1289->1156.
% validate_tex.sh / check_roundtrip.sh OK; compile_pdf.sh OK (102 pages, 0 LaTeX warnings).
%
% Diff (C-57 only; baseline = working tree after R3b-01):
% --- a/sections/introduction.tex
% +++ b/sections/introduction.tex
% @@ -10 +10 @@
% -A critical observation emerges from this body of work: in all published wheeled-legged systems, locomotion mode switching is resolved either heuristically through hand-coded threshold logic \citep{Kashiri2019CENTAURO}, by model predictive control \citep{bjelonic2019,Bjelonic2022OfflineMotion}, or by end-to-end RL policies that encode switching implicitly as a black-box behavior \citep{Lee2024Science,Humphreys2025BioInspired}. None of these approaches employs a neurobiologically interpretable arbitration architecture. Table~\ref{table:action_selection_comparison} summarizes this landscape along axes relevant to the present design choice, horizontally contrasting hand-coded threshold logic, model predictive control, end-to-end reinforcement learning, central pattern generators (CPGs), and basal-ganglia-inspired winner-take-all arbitration. CPGs are predominantly used as a rhythm-\emph{generation} mechanism, but a body of work on feedback-modulated CPGs shows that they can also perform discrete-mode \emph{arbitration} directly, autonomously triggering gait transitions (e.g., walk-to-trot-to-gallop) from vestibular or load feedback without an external switching signal \citep{Fukui2019VestibularCPG,Owaki2017SpontaneousGait} --- the rationale for keeping that role separate from CPG execution in the present design is given below, alongside the CPG description itself.
% +A critical observation emerges from this body of work: in all published wheeled-legged systems, locomotion mode switching is resolved either heuristically through hand-coded threshold logic \citep{Kashiri2019CENTAURO}, by model predictive control \citep{bjelonic2019,Bjelonic2022OfflineMotion}, or by end-to-end RL policies that encode switching implicitly as a black-box behavior \citep{Lee2024Science,Humphreys2025BioInspired}. None of these approaches employs a neurobiologically interpretable arbitration architecture. \revblue{Table~\ref{table:action_selection_comparison} contrasts these approaches with central pattern generators (CPGs) ---including feedback-modulated variants that trigger gait transitions from vestibular or load feedback--- and with the basal-ganglia winner-take-all arbitration adopted here.}
% --- a/sections/methodology.tex
% +++ b/sections/methodology.tex
% @@ -4 +4 @@
% -This study was structured around four main components. First, the robot was designed with a morphology enabling transformation among quadrupedal, differential-driving, and omnidirectional mobile modes. Second, the platform was simulated in the Gazebo environment, integrating the sensors required to perform the assigned tasks. Third, the neuronal control circuits were developed, including modules responsible for obstacle detection via LiDAR, motion-direction definition, and locomotion-mode selection. Additionally, a multilayer perceptron network was implemented to detect terrain instability and adjust the gait mode accordingly. Finally, the physical robot was constructed and implemented following the platform previously designed for simulation, which served as the definitive morphological reference. Manufacturing and assembly insights from the prior engineering development by \citet{Grajales_Des_2025} were incorporated to streamline the fabrication process; however, the simulation design was adopted as the baseline to ensure consistency between the virtual and physical systems. 
% +\revblue{The work comprised four stages: designing a morphology that transforms among quadrupedal, differential-drive, and omnidirectional modes (Section~\ref{ssec:Platform_Design}); simulating it with its sensors in Gazebo (Section~\ref{ssec:Gazebo_Simulation}); developing the neural control circuits for LiDAR obstacle detection, motion direction, and locomotion-mode selection, complemented by a multilayer perceptron (MLP) terrain-instability detector; and building the physical robot (Section~\ref{ssec:Construction}). Fabrication drew on the prior engineering development of \citet{Grajales_Des_2025}, but the simulation design was kept as the morphological baseline to ensure virtual--physical consistency.} 
% @@ -34 +34 @@
% -To evaluate the stability of the platform on irregular terrains, a synthetic terrain was designed featuring multiple inclinations and roughness levels. The platform successfully traversed these surfaces thanks to its mechanical structure and the adaptability of its dynamic gaits. This environment includes three main topologies (Figure~\ref{fig:TERR_IMU}): a first section with flat terrain, a second with gravel-like roughness, and a third consisting of a slope with progressive inclinations of $5^\circ$, $10^{\circ}$, $25^{\circ}$, and $45^{\circ}$ relative to the horizontal. To systematically evaluate the performance of the neural arbitration network, the multi-body physical simulations were strictly restricted to structured topographies, including calibrated inclines, procedurally generated Perlin-noise surfaces, and simple non-dynamic obstacles. This layout ensures a highly controlled and reproducible virtual environment to validate the algorithmic response of the controller before its embedded physical implementation.   
% +To evaluate the stability of the platform on irregular terrains, a synthetic terrain was designed featuring multiple inclinations and roughness levels. The platform successfully traversed these surfaces thanks to its mechanical structure and the adaptability of its dynamic gaits. This environment includes three main topologies (Figure~\ref{fig:TERR_IMU}): a first section with flat terrain, a second with gravel-like roughness, and a third consisting of a slope with progressive inclinations of $5^\circ$, $10^{\circ}$, $25^{\circ}$, and $45^{\circ}$ relative to the horizontal. To systematically evaluate the performance of the neural arbitration network, simulations were restricted to structured topographies, including calibrated inclines, procedurally generated Perlin-noise surfaces, and simple non-dynamic obstacles.   
% @@ -74,3 +74 @@
% -To perceive appetitive and aversive stimuli, a visual input was implemented through a camera, which provided information on localization and, thanks to image processing, allows determining the nature of the presented visual stimulus. 
% -
% -For obstacle detection, a LiDAR sensor was used, as it provides precise information regarding the obstacle’s position. To achieve a sensory transformation of the sensor’s output into neural activations suitable for processing by the corresponding network module, a ring of neurons was implemented to represent the orientation of the detected obstacle. This structure is based on the study by \citet{pardo2022bio}, in which a neural ring encoding the angular orientation of a specific target is proposed. Although inspired by its functionality, the module developed here differs mainly in the Gaussian activations used for the input neurons ---explained later--- and in the inclusion of a decoding structure that determines the intensity of the motor response. The ring consists of 16 units, each responding with maximal amplitude to a preferred angle, indicating the direction in which the obstacle is located.
% +\revblue{As outlined in Section~\ref{ssec:Platform_Design}, the camera identifies and localizes appetitive and aversive stimuli through image processing, while the LiDAR localizes obstacles. To transform the LiDAR output into neural activations,} a ring of neurons was implemented to represent the orientation of the detected obstacle. This structure is based on the study by \citet{pardo2022bio}, in which a neural ring encoding the angular orientation of a specific target is proposed. Although inspired by its functionality, the module developed here differs mainly in the Gaussian activations used for the input neurons ---explained later--- and in the inclusion of a decoding structure that determines the intensity of the motor response. The ring consists of 16 units, each responding with maximal amplitude to a preferred angle, indicating the direction in which the obstacle is located.
% @@ -137 +135 @@
% -The activations of the sensory ring are processed by three integrator layers: input, response, and auxiliary (Figure~\ref{fig:IntLidar}). Cross-connections between the input and response layers encode the direction of the evasive movement, while the backward excitatory projections between the response and auxiliary layers define the intensity of the obstacle-related stimulus. These layers determine the predominance of the evasive action ---information that is stored in the accumulator neuron $L_4$. In the following equations, the subscript $(\cdot)_+$ denotes a rectifying operator (positive-part function): for any scalar $x$, $(x)_+ = \max(x,0)$. This ensures that only nonnegative inputs contribute to the layer dynamics, preventing negative net drive from decreasing the activity through this term. The dynamics of this system are described by differential equations (Equations~\eqref{eq:L0}–\eqref{eq:L4}) shown below. 
% +The activations of the sensory ring are processed by three integrator layers: input, response, and auxiliary (Figure~\ref{fig:IntLidar}). Cross-connections between the input and response layers encode the direction of the evasive movement, while the backward excitatory projections between the response and auxiliary layers define the intensity of the obstacle-related stimulus. These layers determine the predominance of the evasive action ---information that is stored in the accumulator neuron $L_4$. In the following equations, the subscript $(\cdot)_+$ denotes a rectifying operator (positive-part function): for any scalar $x$, $(x)_+ = \max(x,0)$. The dynamics of this system are described by differential equations (Equations~\eqref{eq:L0}–\eqref{eq:L4}) shown below. 
% @@ -150,5 +148 @@
% -In Equations~\eqref{eq:L0}–\eqref{eq:L4} the subscripts are used to enumerate the neurons in each layer of the obstacle sensory module. The variable $Z_i$ denotes the activity of the $i$-th neuron of the LiDAR-associated sensory ring, with $i$ ordered angularly around the robot. The summations $\sum_{i=a}^{b} Z_i$ group contiguous subsets of ring neurons that encode directional sectors (front, back, left, and right); the tank neurons $L_0,\dots,L_3$ integrate these groupings and store the predominance of each directional quadrant (Figure~\ref{fig:lidar-combined}).
% -
% -Likewise, the indices $n$ and $r$ enumerate the neurons in the input, response, and auxiliary integrating layers, so that $I_n$, $R_n$, and $A_n$ represent the activity of the $n$-th neuron in each of these layers. In Equations~\eqref{eq:ResL}–\eqref{eq:AuxL}, each neuron $R_n$ and $A_n$ receives weighted contributions from all other neurons through the synaptic weights $W^{(\mathrm{IR})}_{n,i}$, $W^{(\mathrm{RR})}_{n,j}$, and $W^{(\mathrm{RA})}_{n,j}$, where the indices $i$ and $j$ are used only as summation variables. Finally, in Equation~\eqref{eq:L4} the index $r$ runs over all units $A_n$ to accumulate their activity in neuron $L_4$, which acts as a global integrator of the obstacle-stimulus intensity.
% -
% -The matrices $W^{(IR)}$, $W^{(RR)}$, and $W^{(RA)}$ denote connectivity weight matrices (binary or fixed patterns) representing fixed structural inter-neuron connectivity in the Integrating Neural Network (see Figure~\ref{fig:IntLidar}), fixed and uniform in the same sense as the ring-to-integrator weights above.
% +\revblue{Here $Z_i$ is the $i$-th ring neuron, ordered angularly around the robot; the sums $\sum_{i=a}^{b} Z_i$ group the ring into the front, rear, left, and right sectors integrated by $L_0$--$L_3$ (Figure~\ref{fig:lidar-combined}). $I_n$, $R_n$, and $A_n$ are the $n$-th neurons of the input, response, and auxiliary layers; $i$, $j$, and $r$ are summation indices, with $r$ in Equation~\eqref{eq:L4} running over all auxiliary units. $W^{(IR)}$, $W^{(RR)}$, and $W^{(RA)}$ are fixed (binary or patterned) structural connectivity matrices of the Integrating Neural Network (Figure~\ref{fig:IntLidar}).}
% @@ -168 +162 @@
% -The STN processes incoming signals from the visual centers and determines which stimulus has the greatest relevance when multiple stimuli are present. The GPe inhibits lower-intensity behaviors and reinforces the prioritized action, while the GPi complements this mechanism by suppressing competing responses to ensure that only a single action is executed. This iterative mechanism of inhibition and competition guarantees the effective selection of one type of stimulus. The differential equations~\eqref{eq:stn}–\eqref{eq:str} that describe the dynamics of the STN, GPi, GPe, and striatum (STR) connections, along with the hyperparameters of the structure (see Table~\ref{table:Ganglia}), are presented below. 
% +The STN processes incoming signals from the visual centers and determines which stimulus has the greatest relevance when multiple stimuli are present. The GPe inhibits lower-intensity behaviors and reinforces the prioritized action, while the GPi complements this mechanism by suppressing competing responses to ensure that only a single action is executed. This iterative mechanism of inhibition and competition guarantees the effective selection of one type of stimulus. 
% @@ -222 +216 @@
% -where $\text{L}_2$ and $L_4$ denote the left/right directional quadrant neuron and the LiDAR accumulator neuron from Equations~\eqref{eq:L2} and~\eqref{eq:L4}, respectively, and the bracketed term encodes a $20^\circ$ angular margin defining the frontal zone of the robot's field of view ($\tau$ and $\omega_{L\rightarrow X}$ are defined in Table~\ref{table:Locom_NN}'s notation below).
% +where $\text{L}_2$ and $L_4$ denote the left/right directional quadrant neuron and the LiDAR accumulator neuron from Equations~\eqref{eq:L2} and~\eqref{eq:L4}, respectively ($\tau$ and $\omega_{L\rightarrow X}$ are defined in Table~\ref{table:Locom_NN}'s notation below).
% @@ -233 +227 @@
% -The differential equations~\eqref{eq:Z_3}–\eqref{eq:Z_17} describing the dynamics of the locomotion decision module, as well as the corresponding hyperparameters in Table~\ref{table:Locom_NN}, are presented below. In Equation~\eqref{eq:Z_13}, $HI$ denotes an intrinsic hunting instinct that provides a baseline forward/exploratory motivation, allowing the platform to keep moving and exploring even in the absence of external stimuli.
% +In Equation~\eqref{eq:Z_13}, $HI$ denotes an intrinsic hunting instinct that provides a baseline forward/exploratory motivation, allowing the platform to keep moving and exploring even in the absence of external stimuli.
% @@ -299 +293 @@
% -$HI = 3.0$ defines a baseline forward bias to overcome motor static friction. The Naka--Rushton half-saturation constants $\sigma_M = 0.3$ and $\sigma_I = 1.5$ serve distinct roles: $\sigma_M$ provides rapid mode switching through a narrow saturation range, while $\sigma_I$ affords greater noise tolerance in the IMU pathway. Finally, the MLP thresholds for gait switching were empirically determined based on the observed distribution of $\sigma_{az}$ values during stable versus unstable locomotion, ensuring a clear decision boundary while minimizing false positives. These parameter choices reflect a balance between biological plausibility, computational efficiency, and practical performance requirements for the target application domain.
% +$HI = 3.0$ defines a baseline forward bias to overcome motor static friction. The Naka--Rushton half-saturation constants $\sigma_M = 0.3$ and $\sigma_I = 1.5$ serve distinct roles: $\sigma_M$ provides rapid mode switching through a narrow saturation range, while $\sigma_I$ affords greater noise tolerance in the IMU pathway. Finally, the MLP thresholds for gait switching were empirically determined based on the observed distribution of $\sigma_{az}$ values during stable versus unstable locomotion, ensuring a clear decision boundary while minimizing false positives.
% @@ -314 +308 @@
% -Units $z_0$--$z_2$ and $z_{14}$--$z_{16}$ use a second-order Naka--Rushton nonlinearity to transform their net input into a bounded activation. The general form is given in Equation~\eqref{eq:Naka}, where $(\cdot)_+$ denotes the rectified (positive-part) operator.
% +Units $z_0$--$z_2$ and $z_{14}$--$z_{16}$ use a second-order Naka--Rushton nonlinearity to transform their net input into a bounded activation. The general form is given in Equation~\eqref{eq:Naka}.
% @@ -323,15 +317 @@
% -Based on this general formulation, two task-specific activation functions are defined by instantiating Equation~\eqref{eq:Naka} with different sensitivity parameters. Specifically, $f_{\mathrm{Imu}}(\cdot)$ and $f_{\mathrm{Mode}}(\cdot)$ (Equations~\eqref{eq:NakaImu} and~\eqref{eq:NakaMode}) share the same second-order Naka--Rushton structure, but use different saturation constants ($\sigma_I$ and $\sigma_M$) to match the dynamic range required by each neuronal pathway:
% -
% -\begin{equation}
% -f_{\mathrm{Imu}}(x) = \frac{\mathrm{A}\, x_+^2}{\sigma_{I}^2 + x_+^2}
% -\label{eq:NakaImu}
% -\end{equation}
% -
% -
% -\begin{equation}
% -f_{\mathrm{Mode}}(x) = \frac{\mathrm{A}\, x_+^2}{\sigma_{M}^2 + x_+^2}
% -\label{eq:NakaMode}
% -\end{equation}
% -
% -
% -The differential equations~\eqref{eq:Z0}–\eqref{eq:Z16} describing the dynamics of the gait-mode network, along with the hyperparameters in Table~\ref{table:Locom_NN}, are presented below. 
% +\revblue{Two instantiations of Equation~\eqref{eq:Naka} are used, $f_{\mathrm{Imu}}(x) = f(x;\sigma_I)$ and $f_{\mathrm{Mode}}(x) = f(x;\sigma_M)$, whose different half-saturation constants match the dynamic range of the IMU and mode-selection pathways.} 
% @@ -396 +376 @@
% -Using the new information, the gait-decision network was adjusted (Figure~\ref{fig:Modos_LocomocionV2}). In this version, the connections of units $X_0$–$X_2$ were modified, with unit $X_2$ removed from the design. This unit becomes unnecessary only in the physical implementation, because one output of the MLP identifies whether the robotic platform is tilted in both roll and pitch. Thus, for the physical gait-decision module, input $X_2$ becomes obsolete. This distinction results in separate gait-decision modules for simulation and physical implementation. Inputs $X_0$ and $X_1$ now correspond to outputs $N_0$ and $N_1$, respectively. 
% +\revblue{Using this information, the gait-decision network was adjusted for the physical implementation only (Figure~\ref{fig:Modos_LocomocionV2}): unit $X_2$ is removed, since one MLP output already flags combined roll/pitch tilt, so simulation and hardware use separate gait-decision modules; inputs $X_0$ and $X_1$ now correspond to $N_0$ and $N_1$.} 
% @@ -407 +387 @@
% -The hyperparameters (Table~\ref{table:Locom_MLP}) and the Equations~\eqref{eq:Z0V2}–\eqref{eq:Z16V2} reflecting the changes between the first and the newly proposed network structure are presented below. 
% +\revblue{Equations~\eqref{eq:Z0V2}--\eqref{eq:Z16V2} give the modified units.} 
% @@ -419 +399 @@
% -In Equations~\eqref{eq:Z0V2}–\eqref{eq:Z16V2}: $N_0$ is the MLP output probability of flat terrain; $N_1$ is the MLP output probability of inclined terrain; $N_2$ is the binary stagnation indicator ($1 ={}$stuck, $0 ={}$not stuck); and $U_{N_0}$, $U_{N_1}$ are the corresponding activation thresholds (see Table~\ref{table:Locom_MLP}).
% +\revblue{In Equations~\eqref{eq:Z0V2}--\eqref{eq:Z16V2}, $U_{N_0}$ and $U_{N_1}$ are the activation thresholds of $N_0$ and $N_1$ (Table~\ref{table:Locom_MLP}).}
% @@ -460 +440 @@
% -To address the reviewer's request for a finite state machine (FSM) baseline under identical experimental conditions, a second mode-arbitration controller was implemented that inherits the complete sensory pipeline, navigation dynamics, and noise-injection model described above unchanged, substituting only the mode-selection mechanism: instead of the continuous winner-take-all competition of Section~\ref{sec:methodology}'s Basal Ganglia Module, the FSM selects the active mode via a static, ordered priority hierarchy (terrain instability, then obstacle proximity, then aversive stimulus, then appetitive stimulus, defaulting to differential-rolling exploration), each stimulus-triggered rule holding its selected mode for a fixed persistence interval once triggered. Because every other component is shared byte-for-byte between the two controllers, any behavioral difference between them is attributable exclusively to this mode-selection mechanism, not to differences in perception, low-level navigation, or experimental configuration. Both controllers were evaluated across the same seven simulated suites under matching trial counts, noise levels, and success criteria; results are reported in Section~\ref{sec:results}.
% +\revblue{To provide a finite state machine (FSM) baseline under identical experimental conditions,} a second mode-arbitration controller was implemented that inherits the complete sensory pipeline, navigation dynamics, and noise-injection model described above unchanged, substituting only the mode-selection mechanism: instead of the continuous winner-take-all competition of Section~\ref{sec:methodology}'s Basal Ganglia Module, the FSM selects the active mode via a static, ordered priority hierarchy (terrain instability, then obstacle proximity, then aversive stimulus, then appetitive stimulus, defaulting to differential-rolling exploration), each stimulus-triggered rule holding its selected mode for a fixed persistence interval once triggered. Because every other component is shared byte-for-byte between the two controllers, any behavioral difference between them is attributable exclusively to this mode-selection mechanism, not to differences in perception, low-level navigation, or experimental configuration. Both controllers were evaluated across the same seven simulated suites under matching trial counts, noise levels, and success criteria; results are reported in Section~\ref{sec:results}.
% @@ -553 +533 @@
% -These three measures are supplemented by lower-level safeguards common to servo-actuated platforms: an individually calibrated pulse-width mapping per actuator (compensating unit-to-unit mechanical and mounting variation across the twelve joints), a software saturation limit on every commanded joint angle, and a rejection of inverse-kinematics solutions that fall outside the physically reachable workspace before any command reaches the actuator. Together, this layered strategy --- static stability by gait design, torque-aware trajectory shaping, and feed-forward compliance correction --- constitutes the compensating control approach for the platform's low-torque, open-loop actuation constraints. It does not substitute for closed-loop torque or attitude feedback during active balance recovery on severe terrain, which remains an open limitation of the present design (Section~\ref{ssec:Limitations}).
% +These three measures are supplemented by lower-level safeguards common to servo-actuated platforms: an individually calibrated pulse-width mapping per actuator (compensating unit-to-unit mechanical and mounting variation across the twelve joints), a software saturation limit on every commanded joint angle, and a rejection of inverse-kinematics solutions that fall outside the physically reachable workspace before any command reaches the actuator. \revblue{Together, these measures compensate for the platform's low-torque, open-loop actuation but do not substitute for} closed-loop torque or attitude feedback during active balance recovery on severe terrain, which remains an open limitation of the present design (Section~\ref{ssec:Limitations}).
% --- a/sections/results.tex
% +++ b/sections/results.tex
% @@ -4 +4 @@
% -Initially, simulations of the platform were performed in the Gazebo environment. In this scenario, the expected behaviors of the robotic platform were evaluated under the presentation of different stimuli. Additionally, gait-switching tests on complex terrain were conducted using the three-topology map described earlier. Subsequently, the behaviors observed in simulation were validated using the physical robotic platform constructed as the digital twin counterpart. 
% +\revblue{Results are reported first for Gazebo simulation (single and multiple stimuli, and gait switching on the three-topology terrain map described earlier) and then for the physical platform.} 
% @@ -82 +82 @@
% -\textbf{Aversive Escape (see Figure~\ref{fig:aversive_micro}):} The defensive repulsion properties of the architecture were re-evaluated under the combined-perturbation sweep, with every retained trial across all five severity levels resolving the threat successfully (Figure~\ref{fig:noise_robustness_grid}). The neural firing variance exhibits a swift, step-like onset upon threat detection — structural mobilization into a defensive overdrive state — and returns to baseline once the threat leaves the field of view, demonstrating the efficiency of the neural gating mechanism.
% +\textbf{Aversive Escape (see Figure~\ref{fig:aversive_micro}):} The defensive repulsion properties of the architecture were re-evaluated under the combined-perturbation sweep, with every retained trial across all five severity levels resolving the threat successfully (Figure~\ref{fig:noise_robustness_grid}). The neural firing variance exhibits a swift, step-like onset upon threat detection — structural mobilization into a defensive overdrive state — and returns to baseline once the threat leaves the field of view.
% @@ -91 +91 @@
% -\textbf{Obstacle Evasion (see Figure~\ref{fig:obstacle_micro}):} Smooth capsule geometries for simulated obstacles prevent localized torque spikes during contact. Every retained trial across all five combined-perturbation levels completed evasion successfully (Figure~\ref{fig:noise_robustness_grid}). The sensory-locomotor coupling is precise: neural firing variance stays at a low baseline until the robot enters the obstacle's ranging threshold, at which point the avoidance sub-circuits modulate steering weights while preserving forward momentum.
% +\textbf{Obstacle Evasion (see Figure~\ref{fig:obstacle_micro}):} Smooth capsule geometries for simulated obstacles prevent localized torque spikes during contact. Every retained trial across all five combined-perturbation levels completed evasion successfully (Figure~\ref{fig:noise_robustness_grid}). Neural firing variance stays at a low baseline until the robot enters the obstacle's ranging threshold, at which point the avoidance sub-circuits modulate steering weights while preserving forward momentum.
% @@ -128 +128 @@
% -\textbf{Complex Navigation (see Figure~\ref{fig:noise_robustness_grid}):} The complex navigation suite culminates the simulation validation, presenting the obstacle, aversive, and appetitive stimuli within a single trial. In the neural firing-variance profile of Figure~\ref{fig:complex_micro}, the robot first clears the physical obstacle through the LiDAR-driven evasion pathway, without a corresponding shaded region since this channel is not one of the two stimulus-presence traces plotted; the aversive threat then drives a red-shaded burst of firing-variance activity as the quadrupedal escape response engages; once clear of the threat, a blue-shaded interval marks sustained appetitive tracking as the platform reorients and approaches the target under differential-rolling mode. The kinematic panel's Roll/Pitch RMS response tracks this same three-stimulus sequence, rising through each gait transition and settling once the target is reached. Every retained trial across all five combined-perturbation severity levels resolves this three-stimulus encounter successfully (Figure~\ref{fig:noise_robustness_grid}), evidencing that the arbitration network does not suffer from computational deadlocks even when a single trial exercises all three behavioral channels under joint sensor perturbation.
% +\textbf{Complex Navigation (see Figure~\ref{fig:noise_robustness_grid}):} The complex navigation suite presents the obstacle, aversive, and appetitive stimuli within a single trial. In the neural firing-variance profile of Figure~\ref{fig:complex_micro}, the robot first clears the physical obstacle through the LiDAR-driven evasion pathway, without a corresponding shaded region since this channel is not one of the two stimulus-presence traces plotted; the aversive threat then drives a red-shaded burst of firing-variance activity as the quadrupedal escape response engages; once clear of the threat, a blue-shaded interval marks sustained appetitive tracking as the platform reorients and approaches the target under differential-rolling mode. The kinematic panel's Roll/Pitch RMS response tracks this same three-stimulus sequence, rising through each gait transition and settling once the target is reached. Every retained trial across all five combined-perturbation severity levels resolves this three-stimulus encounter successfully (Figure~\ref{fig:noise_robustness_grid}), evidencing that the arbitration network does not suffer from computational deadlocks even when a single trial exercises all three behavioral channels under joint sensor perturbation.
% @@ -221 +221 @@
% -For the experiments involving variable and unstable topographies in simulation, a three‐section terrain map was designed, as shown in Figure~\ref{fig:TERR_IMU}. In the first simulated experiment, the robot begins moving forward driven by the hunting instinct encoded in the locomotor decision network. Subsequently, the transition from regular to rough terrain is introduced to analyze the standard deviation of the acceleration along the $Z$ axis throughout the simulation. In the physical robot, this information corresponds to the MLP outputs: $N_{0}$ (flat-terrain probability), $N_{1}$ (inclined-terrain probability), and $N_{2}$ (binary stagnation indicator). 
% +For the experiments involving variable and unstable topographies in simulation, a three‐section terrain map was designed, as shown in Figure~\ref{fig:TERR_IMU}. In the first simulated experiment, the robot begins moving forward driven by the hunting instinct encoded in the locomotor decision network. Subsequently, the transition from regular to rough terrain is introduced to analyze the standard deviation of the acceleration along the $Z$ axis throughout the simulation. \revblue{In the physical robot, this information corresponds to the MLP outputs $N_0$--$N_2$ (Section~\ref{ssec:Gait_Decision_Module}).} 
% @@ -315 +315 @@
% -To quantitatively evaluate the dynamic stability and structural safety of the hybrid platform during high-transient morphological changes, the Tipover Risk ($TR$) index was tracked and unified across 15 randomized experimental runs for both rugged terrain and inclined slope, as depicted in Figure~\ref{fig:stability_profile}. The $TR$ metric was formally derived from the McGhee--Frank Normalized Stability Margin as $TR = 1 - SM/r_{in}$, where $SM$ is the shortest horizontal distance between the center of mass projection and the support polygon boundary, and $r_{in}$ is the inradius of that polygon. Values approaching $TR = 0.0$ represent maximum structural safety, while $TR = 1.0$ marks the critical overturning threshold. By aligning the temporal records around the exact neural mode-switch instant ($t = 0.0\,\text{s}$), the physical efficacy of the bio-inspired controller was made directly observable. In both terrains, the platform maintained a consistently bounded risk profile around $\mu_{TR} \approx 0.58$--$0.60$ (rugged terrain mean $0.600$; inclined slope mean $0.593$), well below the critical overturning threshold throughout the recorded window. Rather than a monotonic escalation, the rugged-terrain trace showed an isolated transient spike reaching $TR \approx 0.77$ shortly before the mode-switch instant ($t \approx -1.3$ to $-0.5\,\text{s}$), consistent with the periodic contraction of the support polygon as individual legs were lifted during the quadrupedal gait cycle, before settling back to the $\approx 0.59$ baseline across the transition at $t = 0.0\,\text{s}$. The inclined-slope trace remained flat at $TR \approx 0.59$ across the entire window, with no escalation or drop coincident with the mode switch. Across all 30 trials, the $\pm 1\sigma$ statistical envelope never crossed the critical overturning boundary ($TR = 1.0$), providing empirical validation of the safety and repeatability of the control architecture under stochastic noise conditions.
% +\revblue{The Tipover Risk ($TR$) index was tracked across 15 randomized runs per terrain (Figure~\ref{fig:stability_profile}).} The $TR$ metric was formally derived from the McGhee--Frank Normalized Stability Margin as $TR = 1 - SM/r_{in}$, where $SM$ is the shortest horizontal distance between the center of mass projection and the support polygon boundary, and $r_{in}$ is the inradius of that polygon. Values approaching $TR = 0.0$ represent maximum structural safety, while $TR = 1.0$ marks the critical overturning threshold. In both terrains, the platform maintained a consistently bounded risk profile around $\mu_{TR} \approx 0.58$--$0.60$ (rugged terrain mean $0.600$; inclined slope mean $0.593$) throughout the recorded window. Rather than a monotonic escalation, the rugged-terrain trace showed an isolated transient spike reaching $TR \approx 0.77$ shortly before the mode-switch instant ($t \approx -1.3$ to $-0.5\,\text{s}$), consistent with the periodic contraction of the support polygon as individual legs were lifted during the quadrupedal gait cycle, before settling back to the $\approx 0.59$ baseline across the transition at $t = 0.0\,\text{s}$. The inclined-slope trace remained flat at $TR \approx 0.59$ across the entire window, with no escalation or drop coincident with the mode switch. \revblue{Across all 30 trials the $\pm 1\sigma$ envelope never crossed $TR = 1.0$.}
% @@ -320 +320 @@
% -To quantitatively assess the adaptive morphology and environmental specialization of the control architecture, the steady-state kinematic behaviors of both experimental terrains were mapped onto a phase space defined by the chassis Roll RMS ($\phi_{\text{RMS}}$) and Pitch RMS ($\theta_{\text{RMS}}$) across all 30 successful trials, as shown in Figure~\ref{fig:terrain_adaptability}. The results revealed a clear separation in Pitch RMS between the two terrains, with partial overlap in Roll RMS. The rugged terrain converged into a low-to-moderate roll, low-pitch operational zone ($\mu_{\phi} \approx 1.59^{\circ}$, $\mu_{\theta} \approx 1.67^{\circ}$), capturing the kinematic footprint of asymmetric foot-ground interactions on uneven stones while maintaining a largely level chassis. Conversely, the inclined slope clustered in a higher-pitch region with comparable roll magnitude ($\mu_{\phi} \approx 1.47^{\circ}$, $\mu_{\theta} \approx 2.29^{\circ}$). This spatial distribution reflects the geometric adaptation triggered by the Hybrid morphological transformation: during slope ascent, the chassis aligned longitudinally with the ramp gradient (higher Pitch RMS), while the deployment of the coaxial wheel base kept the lateral support polygon wide enough that roll disturbances remained similar in magnitude to the rugged-terrain baseline rather than being further suppressed. The Pitch-RMS separation between both domains provides empirical evidence that the bio-inspired network selected a distinct, repeatable, terrain-specialized kinematic posture for each environmental condition, rather than applying a generic, environment-agnostic control law, while the Roll-RMS overlap indicates that lateral stabilization was comparably effective across both terrains rather than being terrain-specific.
% +\revblue{The steady-state kinematics of both terrains were mapped onto a Roll RMS ($\phi_{\text{RMS}}$) vs.\ Pitch RMS ($\theta_{\text{RMS}}$) phase space across all 30 successful trials (Figure~\ref{fig:terrain_adaptability}).} The results revealed a clear separation in Pitch RMS between the two terrains, with partial overlap in Roll RMS. The rugged terrain converged into a low-to-moderate roll, low-pitch operational zone ($\mu_{\phi} \approx 1.59^{\circ}$, $\mu_{\theta} \approx 1.67^{\circ}$), capturing the kinematic footprint of asymmetric foot-ground interactions on uneven stones while maintaining a largely level chassis. Conversely, the inclined slope clustered in a higher-pitch region with comparable roll magnitude ($\mu_{\phi} \approx 1.47^{\circ}$, $\mu_{\theta} \approx 2.29^{\circ}$). This spatial distribution reflects the geometric adaptation triggered by the Hybrid morphological transformation: during slope ascent, the chassis aligned longitudinally with the ramp gradient (higher Pitch RMS), while the deployment of the coaxial wheel base kept the lateral support polygon wide enough that roll disturbances remained similar in magnitude to the rugged-terrain baseline rather than being further suppressed. \revblue{The Pitch-RMS separation indicates that the network selects a distinct, repeatable posture for each terrain rather than applying an environment-agnostic control law.}
% @@ -326,13 +326 @@
% -To address the reviewer's request for systematic quantitative metrics
% -on mode-switching behavior, a dedicated stability analysis was
% -conducted covering the Quadrupedal Walking gait and all six pairwise transitions
% -among the three locomotion modes.
% -
% -\paragraph{Stability metrics.}
% -
% -The Normalized Stability Margin ($SM$) and Tipover Risk
% -($TR = 1 - SM/r_{in}$) are defined in
% -Section~\ref{sssec:variable_topo_quant}. Here, the same metrics were
% -derived frame-by-frame from the foot-contact positions logged during
% -simulation for the Quadrupedal Walking gait and all six pairwise
% -mode transitions.
% +\revblue{This subsubsection quantifies static stability for the Quadrupedal Walking gait and all six pairwise transitions among the three locomotion modes, using the Normalized Stability Margin ($SM$) and Tipover Risk ($TR$) of Section~\ref{sssec:variable_topo_quant}, computed frame-by-frame from the foot-contact positions logged in simulation.}
% @@ -422,7 +410 @@
% -Figure~\ref{fig:transition_phases} shows the three undirected mode pairs
% -($C \leftrightarrow H$, $H \leftrightarrow X$, $X \leftrightarrow C$)
% -side by side, each with a double-headed arrow denoting that the
% -transition is reversible: the two directions between a given pair of
% -modes (e.g. $C \to H$ and $H \to C$) converge to statistically
% -indistinguishable support-polygon geometries and Tipover Risk values,
% -so each pair is shown once rather than as two separate directed rows.
% +\revblue{Both directions of each mode pair converge to statistically indistinguishable support-polygon geometries and $TR$ values, so Figure~\ref{fig:transition_phases} shows each undirected pair ($C \leftrightarrow H$, $H \leftrightarrow X$, $X \leftrightarrow C$) once.}
% @@ -461,2 +442,0 @@
% -The support polygon therefore remains non-degenerate throughout the
% -entire switching interval.
% @@ -465 +445 @@
% -$N_{\mathrm{total}} = 979$ logged samples, every transition completed
% +logged samples, every transition completed
% @@ -479,5 +459 @@
% -Regarding oscillation amplitude during transitions:
% -because all morphological reconfiguration is performed with the
% -full foot set in ground contact, no abrupt postural discontinuity
% -is introduced at the switching instant, regardless of the chassis
% -velocity at that moment.
% +\revblue{Because reconfiguration is performed with all four feet in contact, no abrupt postural discontinuity is introduced at the switching instant.}
% @@ -496,10 +472 @@
% -Regarding energy consumption:
% -the Gazebo simulation environment was not instrumented with a
% -power-monitoring middleware plugin; consequently, joint-level energy
% -data cannot be extracted from the virtual trials reported in this
% -subsubsection.
% -Energy consumption across locomotion modes is, however, addressed
% -in the physical implementation section of the present work, where
% -per-motor current and power measurements were acquired during the
% -physical platform experiments; the reader is referred to that
% -section for the corresponding energy characterization of each mode.
% +\revblue{Joint-level energy was not logged in Gazebo (no power-monitoring plugin); energy per mode is characterized on the physical platform in Section~\ref{ssec:QuantReal}.}
% @@ -511 +478 @@
% -To address the reviewer's request for a finite state machine (FSM) baseline under identical experimental conditions, the priority-hierarchy controller described in Section~\ref{ssec:FSMBaseline} was evaluated against the basal-ganglia arbitration network across all seven simulated suites, matched on trial count, noise levels, and success criteria. Table~\ref{tab:fsm_comparison} reports success rate, task completion time, Roll RMS, and $\lambda$ --- a motor-command clarity ratio (commanded-velocity magnitude relative to neural activity norm; higher indicates more decisive output) --- for both controllers in each suite.
% +\revblue{Table~\ref{tab:fsm_comparison} compares the FSM controller of Section~\ref{ssec:FSMBaseline} with the basal-ganglia network on all seven simulated suites, reporting success rate, task completion time, Roll RMS, and $\lambda$ --- a motor-command clarity ratio (commanded-velocity magnitude relative to neural activity norm; higher indicates more decisive output).}
% @@ -539 +506 @@
% -In the three single-stimulus suites (Appetitive, Aversive, Obstacle), the FSM's fixed priority rule is structurally well-matched to the task --- one stimulus is active at a time, so a static hierarchy resolves it immediately and deterministically. Both controllers eventually complete every mission in all three suites (100\% mission-level success in Table~\ref{tab:fsm_comparison}, i.e.\ a trial that first times out and is then retried still counts as a success once it lands, and this is the only sense in which "success rate" is used anywhere in this paper). The two controllers diverge, however, in how often that retry is needed: the FSM never requires one, at any noise level, while the Neural system occasionally does as noise increases. Figure~\ref{fig:attempts_vs_noise_level} makes this explicit with its own metric, attempts needed per completed mission (not a success rate) --- Neural needs up to $1.67\times$ as many attempts on average in Appetitive and $1.33\times$ in Aversive at the noisiest levels tested, non-monotonically given the small per-level sample, while FSM stays at exactly one attempt throughout.
% +In the three single-stimulus suites (Appetitive, Aversive, Obstacle), the FSM's fixed priority rule is structurally well-matched to the task --- one stimulus is active at a time, so a static hierarchy resolves it immediately and deterministically. Both controllers eventually complete every mission in all three suites (100\% mission-level success in Table~\ref{tab:fsm_comparison}, i.e.\ a trial that first times out and is then retried still counts as a success once it lands, and this is the only sense in which "success rate" is used anywhere in this paper). \revblue{The two controllers diverge, however, in how often that retry is needed (Figure~\ref{fig:attempts_vs_noise_level}; attempts needed per completed mission, not a success rate):} Neural needs up to $1.67\times$ as many attempts on average in Appetitive and $1.33\times$ in Aversive at the noisiest levels tested, non-monotonically given the small per-level sample, while FSM stays at exactly one attempt throughout.
% @@ -566 +533 @@
% -To directly address the reviewer's request for inspection-oriented task indicators, \revblue{the platform was evaluated in the dedicated obstacle-field scenario described in Section~\ref{ssec:Gazebo_Simulation}, in which the target is not visible from the initial pose (sampling rule in Table~\ref{tab:inspection_metrics}).}
% +\revblue{To obtain inspection-oriented task indicators, the platform was evaluated in the dedicated obstacle-field scenario described in Section~\ref{ssec:Gazebo_Simulation}, in which the target is not visible from the initial pose (sampling rule in Table~\ref{tab:inspection_metrics}).}
% @@ -568 +535 @@
% -As described in Section~\ref{ssec:Gazebo_Simulation}, the system's behavior in this task is purely reactive --- no map, no spatial memory, no deliberate re-routing when blocked. What it implements instead is reactive re-routing driven by the basal-ganglia/winner-take-all (BG/WTA) arbitration network, which integrates LiDAR and camera information in real time to laterally evade each obstacle while continuing toward the target once visible. Coverage rate --- the percentage of a predefined inspection area visited --- is not reported here, as it presupposes a coverage planner and a known target area that this reactive architecture neither has nor requires (see Section~\ref{ssec:Limitations}). Instead, task-completion indicators that are both measurable with the existing instrumentation and consistent with the system's actual capabilities are used: success rate (SR), final approach precision ($d_{final}$), and time to first visual acquisition of the target ($T_{acquisition}$).
% +\revblue{As described in Section~\ref{ssec:Gazebo_Simulation}, behavior in this task is purely reactive (no map, no spatial memory); re-routing emerges from the basal-ganglia/winner-take-all (BG/WTA) arbitration network integrating LiDAR and camera input in real time.} Coverage rate --- the percentage of a predefined inspection area visited --- is not reported here, as it presupposes a coverage planner and a known target area that this reactive architecture neither has nor requires. Instead, task-completion indicators that are both measurable with the existing instrumentation and consistent with the system's actual capabilities are used: success rate (SR), final approach precision ($d_{final}$), and time to first visual acquisition of the target ($T_{acquisition}$).
% @@ -572 +539 @@
% -Twenty trials yielded SR $= 85\%$ (17/20), with three timeouts (trials 7, 11, 17) discussed below. Among the 17 successful trials, $d_{final} = 1.076 \pm 0.038$~m (Table~\ref{tab:inspection_metrics}), directly addressing the reviewer's request for target tracking accuracy: the BG$\rightarrow X_{17}$ loop converges to essentially the same stopping distance from the target regardless of the route taken to reach it.
% +Twenty trials yielded SR $= 85\%$ (17/20), with three timeouts (trials 7, 11, 17) discussed below. Among the 17 successful trials, $d_{final} = 1.076 \pm 0.038$~m (Table~\ref{tab:inspection_metrics}): the BG$\rightarrow X_{17}$ loop converges to essentially the same stopping distance from the target regardless of the route taken to reach it.
% @@ -593 +560 @@
% -The 17 successful trials separate into two navigation patterns distinguished by $N_{ev}^{LiDAR}$: a clean pattern ($n=12$, $N_{ev}^{LiDAR} = 4$--$7$, close to one detection per capsule, $d_{final} = 1.081 \pm 0.042$~m) and an elevated-evasion pattern ($n=5$, $N_{ev}^{LiDAR} = 11$--$28$, additional evasive maneuvering before finding a route toward the target, $d_{final} = 1.065 \pm 0.022$~m). Despite the substantially messier routes, the elevated-evasion group converges at least as tightly on the stopping point as the clean group (Figure~\ref{fig:inspection_dfinal_scatter}) --- if anything, more tightly, though the small elevated-group sample ($n=5$) cautions against over-reading that difference. What both subgroups clearly share is that route complexity does not degrade final approach precision, confirming the stopping criterion is robust across the observed range of $N_{ev}^{LiDAR}$. The three timeouts are extreme cases of the elevated pattern ($N_{ev}^{LiDAR} = 18$, $27$, $33$): in all three, the target was visually detected ($T_{acquisition}$ non-null) but the appetitive channel never accumulated sufficient sustained activation to trigger $X_{17}$ before the deadline, as competing obstacle-channel activation persisted throughout. This is consistent with the architecture's known limitation --- reactive navigation without spatial memory can fail under unfavorable dynamic configurations (Section~\ref{ssec:Limitations}) --- rather than a failure of the inspection-stop mechanism itself.
% +The 17 successful trials separate into two navigation patterns distinguished by $N_{ev}^{LiDAR}$: a clean pattern ($n=12$, $N_{ev}^{LiDAR} = 4$--$7$, close to one detection per capsule, $d_{final} = 1.081 \pm 0.042$~m) and an elevated-evasion pattern ($n=5$, $N_{ev}^{LiDAR} = 11$--$28$, additional evasive maneuvering before finding a route toward the target, $d_{final} = 1.065 \pm 0.022$~m). Despite the substantially messier routes, the elevated-evasion group converges at least as tightly on the stopping point as the clean group (Figure~\ref{fig:inspection_dfinal_scatter}) --- if anything, more tightly, though the small elevated-group sample ($n=5$) cautions against over-reading that difference. The three timeouts are extreme cases of the elevated pattern ($N_{ev}^{LiDAR} = 18$, $27$, $33$): in all three, the target was visually detected ($T_{acquisition}$ non-null) but the appetitive channel never accumulated sufficient sustained activation to trigger $X_{17}$ before the deadline, as competing obstacle-channel activation persisted throughout. This is consistent with the architecture's known limitation --- reactive navigation without spatial memory can fail under unfavorable dynamic configurations --- rather than a failure of the inspection-stop mechanism itself.
% @@ -602 +569 @@
% -This experiment also provides quantitative evidence responding to the reviewer's request for autonomous re-planning after obstacle blockage: \revblue{behavioral mode changes ($N_{ev}^{X}$) were $\approx 2.3\times$ fewer than the sensory-level obstacle detections ($N_{ev}^{LiDAR}$) that triggered them (Table~\ref{tab:inspection_metrics})}, in every one of the 17 successful trials without exception. This demonstrates that the WTA arbitration converts repeated raw sensory fluctuations into a smaller number of sustained behavioral commitments, avoiding the mode-dithering a simple threshold-triggered controller would exhibit --- the mechanism underlying the reactive re-routing observed in this task.
% +\revblue{This experiment also quantifies re-routing after obstacle blockage:} \revblue{behavioral mode changes ($N_{ev}^{X}$) were $\approx 2.3\times$ fewer than the sensory-level obstacle detections ($N_{ev}^{LiDAR}$) that triggered them (Table~\ref{tab:inspection_metrics})}, in every one of the 17 successful trials without exception. This demonstrates that the WTA arbitration converts repeated raw sensory fluctuations into a smaller number of sustained behavioral commitments, avoiding the mode-dithering a simple threshold-triggered controller would exhibit --- the mechanism underlying the reactive re-routing observed in this task.
% @@ -604 +571 @@
% -\revblue{$T_{acquisition}$ (Table~\ref{tab:inspection_metrics})} is a first-acquisition metric, not re-acquisition, since the target is not visible from the initial pose in this scenario (Section~\ref{ssec:Limitations}). Visual pursuit is itself conditioned on the absence of a frontal obstacle: during evasion it is suspended, resuming automatically once the obstacle clears --- the system's re-acquisition-like behavior. The robot remained in the wheeled/differential locomotion mode (H) throughout all 20 trials, so no gait-transition stability data applies to this experiment.
% +\revblue{$T_{acquisition}$ (Table~\ref{tab:inspection_metrics})} is a first-acquisition metric, not re-acquisition, since the target is not visible from the initial pose in this scenario. Visual pursuit is itself conditioned on the absence of a frontal obstacle: during evasion it is suspended, resuming automatically once the obstacle clears --- the system's re-acquisition-like behavior. The robot remained in the wheeled/differential locomotion mode (H) throughout all 20 trials, so no gait-transition stability data applies to this experiment.
% @@ -699 +666 @@
% -As illustrated in Figure~\ref{fig:mlp_eval_grid}a, the final model produced a confusion matrix on the held-out test set. Across the 5-fold cross-validation, the model achieved an accuracy of $0.767 \pm 0.028$, precision of $0.764 \pm 0.024$, recall of $0.779 \pm 0.048$, and an F1-score of $0.771 \pm 0.033$. On the held-out test set ($n = 120$), the model reached an accuracy of 0.800, precision of 0.747, recall of 0.918, F1-score of 0.824, and an ROC-AUC of 0.864, confirming that the network reliably discriminates between flat and inclined terrain, with a recall bias toward the inclined-terrain class that minimizes the safety-critical risk of missed inclines. The classifier's discrimination capacity is further supported by the ROC curve (Figure~\ref{fig:mlp_eval_grid}b), where the curve remains well above the chance diagonal across the operating range, yielding an AUC of 0.864.
% +Across the 5-fold cross-validation, the model achieved an accuracy of $0.767 \pm 0.028$, precision of $0.764 \pm 0.024$, recall of $0.779 \pm 0.048$, and an F1-score of $0.771 \pm 0.033$. On the held-out test set ($n = 120$\revblue{; confusion matrix in Figure~\ref{fig:mlp_eval_grid}a}), the model reached an accuracy of 0.800, precision of 0.747, recall of 0.918, F1-score of 0.824, and an ROC-AUC of 0.864, confirming that the network reliably discriminates between flat and inclined terrain, with a recall bias toward the inclined-terrain class that minimizes the safety-critical risk of missed inclines. \revblue{The ROC curve (Figure~\ref{fig:mlp_eval_grid}b) stays well above the chance diagonal across the operating range.}
% @@ -734 +701 @@
% -The MLP was further benchmarked against alternative classifiers (Logistic Regression, SVM with RBF kernel, and Random Forest) under the same cross-validation protocol. Logistic Regression underperformed ($\text{F1} = 0.692$), reflecting the non-linear separability of the IMU features, while SVM-RBF ($\text{F1} = 0.857$) and Random Forest ($\text{F1} = 0.947$) achieved higher raw scores (Figure~\ref{fig:classifier_tradeoff}a); the rationale for nonetheless selecting the MLP is given below, alongside the per-sample latency comparison.
% +The MLP was further benchmarked against alternative classifiers (Logistic Regression, SVM with RBF kernel, and Random Forest) under the same cross-validation protocol. Logistic Regression underperformed ($\text{F1} = 0.692$), reflecting the non-linear separability of the IMU features, while SVM-RBF ($\text{F1} = 0.857$) and Random Forest ($\text{F1} = 0.947$) achieved higher raw scores (Figure~\ref{fig:classifier_tradeoff}a).
% @@ -736 +703 @@
% -To assess whether the accuracy gap with respect to SVM-RBF and Random Forest justifies their adoption in an embedded, real-time control loop, a computational-cost analysis was conducted measuring per-sample inference latency for each classifier. As shown in Figure~\ref{fig:classifier_tradeoff}b, the MLP requires only $1.64\,\mu\text{s}$ per sample, approximately $14\times$ faster than the SVM-RBF ($23.5\,\mu\text{s}$) and $25\times$ faster than the Random Forest ($41.0\,\mu\text{s}$), while substantially outperforming Logistic Regression ($\text{F1} = 0.692$) at a comparable latency of $0.54\,\mu\text{s}$. \revblue{Given the real-time constraints of the neural-arbitration pipeline and the prospect of deploying the classifier on a resource-constrained embedded microcontroller,} the MLP (150–80) was therefore selected as the configuration offering the best accuracy-to-latency trade-off for terrain classification, rather than the configuration with the highest raw classification score. Figure~\ref{fig:classifier_tradeoff}c additionally reports model size (parameters/tree nodes) for the same four classifiers, a directly relevant factor for deployment on resource-constrained embedded hardware. The $1.64\,\mu\text{s}$ figure above was measured on a host PC (PyTorch)\revblue{, the external laptop described in Section~\ref{ssec:HardwareSpec}, which also ran the neural circuits during the physical experiments}; to directly address the reviewer's request for an onboard deployment analysis, the same forward pass was additionally re-implemented in C and benchmarked on the target ESP32-S3 microcontroller ($N=1000$ repeated inferences), yielding a mean on-device latency of $8.81\,\text{ms}$ ($\text{SD}=2.67\,\mu\text{s}$) and a memory footprint of 24{,}140~B RAM (7.4\% of 320~KB) and 397{,}829~B flash (30.4\% of the board's 4~MB partition, almost entirely the $\sim$47k-parameter model itself). The on-device number is markedly higher than the host figure, consistent with a straightforward, non-vectorized scalar implementation reading weights from flash rather than from a host CPU's cache hierarchy.
% +\revblue{Per-sample inference latency (Figure~\ref{fig:classifier_tradeoff}b) explains why the MLP was nonetheless selected: it requires only} $1.64\,\mu\text{s}$ per sample, approximately $14\times$ faster than the SVM-RBF ($23.5\,\mu\text{s}$) and $25\times$ faster than the Random Forest ($41.0\,\mu\text{s}$)\revblue{, while Logistic Regression, at a comparable $0.54\,\mu\text{s}$, is markedly less accurate}. \revblue{Given the real-time constraints of the neural-arbitration pipeline and the prospect of deploying the classifier on a resource-constrained embedded microcontroller,} the MLP (150–80) was therefore selected as the configuration offering the best accuracy-to-latency trade-off for terrain classification, rather than the configuration with the highest raw classification score. Figure~\ref{fig:classifier_tradeoff}c additionally reports model size (parameters/tree nodes) for the same four classifiers\revblue{, relevant for embedded deployment}. The $1.64\,\mu\text{s}$ figure above was measured on a host PC (PyTorch)\revblue{, the external laptop described in Section~\ref{ssec:HardwareSpec}}; the same forward pass was additionally re-implemented in C and benchmarked on the target ESP32-S3 microcontroller ($N=1000$ repeated inferences), yielding a mean on-device latency of $8.81\,\text{ms}$ ($\text{SD}=2.67\,\mu\text{s}$) and a memory footprint of 24{,}140~B RAM (7.4\% of 320~KB) and 397{,}829~B flash (30.4\% of the board's 4~MB partition, almost entirely the $\sim$47k-parameter model itself). The on-device number is markedly higher than the host figure, consistent with a straightforward, non-vectorized scalar implementation reading weights from flash rather than from a host CPU's cache hierarchy.
% @@ -748 +715 @@
% -In the real implementation (Figure~\ref{fig:RES_IMU_real}), images were captured every 21 seconds. Initially, the robot advances over a region with minimal irregularities ---less than 1.2cm--- which it overcomes during the first 13 seconds. At this point, a stagnation condition is detected that persists for the next 10 seconds. When this stagnation signal reaches its maximum sustained amplitude, the transition to the quadrupedal mode is triggered. In this state, the robot traverses the rough terrain and, around second 86, when the activity of unit $X_{0}$ ceases (Figure~\ref{fig:NEU_IMU_real}), the platform transitions back to the differential rolling mode. In Figure~\ref{fig:GRAPH_IMU_real}, the activations of the MLP outputs and the robot’s stagnation state are shown.
% +In the real implementation (Figure~\ref{fig:RES_IMU_real}), images were captured every 21 seconds. Initially, the robot advances over a region with minimal irregularities ---less than 1.2cm--- which it overcomes during the first 13 seconds. At this point, a stagnation condition is detected that persists for the next 10 seconds. When this stagnation signal reaches its maximum sustained amplitude, the transition to the quadrupedal mode is triggered. In this state, the robot traverses the rough terrain and, around second 86, when the activity of unit $X_{0}$ ceases (Figure~\ref{fig:NEU_IMU_real}), the platform transitions back to the differential rolling mode.
% @@ -788 +755 @@
% -A similar behavior is observed in the experiment with the physical robot (Figure~\ref{fig:RES_IMU2_real}). However, in the real robot's gait-decision network, the inclined terrain is detected before second 46 (Figure~\ref{fig:NEU_IMU2_real}, Locomotion Module panel); consequently, a transient response emerges, increasing the activity of neuron $X_{14}$, which commands the robot to climb the ramp using the differential rolling mode; this mode remains active even after clearing the slope. In Figure~\ref{fig:GRAPH_IMU2_real}, the activations of the MLP outputs and the robot's stagnation state are shown.
% +A similar behavior is observed in the experiment with the physical robot (Figure~\ref{fig:RES_IMU2_real}). However, in the real robot's gait-decision network, the inclined terrain is detected before second 46 (Figure~\ref{fig:NEU_IMU2_real}, Locomotion Module panel); consequently, a transient response emerges, increasing the activity of neuron $X_{14}$, which commands the robot to climb the ramp using the differential rolling mode; this mode remains active even after clearing the slope. \revblue{MLP outputs and stagnation state for both scenarios are shown in Figures~\ref{fig:GRAPH_IMU_real} and~\ref{fig:GRAPH_IMU2_real}.}
% @@ -833 +800 @@
% -To further consolidate the physical validation under complex, multi-section topographies, a dedicated testing suite of 15 continuous repetitive trials was conducted over an experimental track whose structural layout systematically combines four topographic segments: a double-ascent ramp, a horizontal elevated plane, a controlled descent ramp, and a rough terrain section characterized by stochastic, non-uniform height variations generated via Perlin noise. The sequential adaptation of the locomotion arbitration network throughout this compound course is documented in Figure~\ref{fig:real_variable_topography_sequence}, which maps the chronological gait-switching decisions enforced by the bio-inspired basal-ganglia controller across six representative frames.
% +To further consolidate the physical validation under complex, multi-section topographies, a dedicated testing suite of 15 continuous repetitive trials was conducted over an experimental track whose structural layout systematically combines four topographic segments: a double-ascent ramp, a horizontal elevated plane, a controlled descent ramp, and a rough terrain section characterized by stochastic, non-uniform height variations generated via Perlin noise. \revblue{Figure~\ref{fig:real_variable_topography_sequence} and Table~\ref{tab:terrain_sequence} document the resulting gait-switching sequence.}
% @@ -853 +820 @@
% -Across the compound course (double-ascent ramp $\rightarrow$ elevated plane $\rightarrow$ descent ramp $\rightarrow$ Perlin-noise rough terrain, 15 trials), the arbitration network autonomously transitions through the gait sequence summarised in Table~\ref{tab:terrain_sequence} (Figure~\ref{fig:real_variable_topography_sequence}). The critical event occurs in Frames 5--6: rolling-mode wheel slip on the stochastic rough segment triggers a sustained stagnation signal, resolved by reactively redeploying the quadrupedal gait, which restores the support polygon and completes traversal. This closed-loop reactive switching was observed across all 15 trials, demonstrating reliable, autonomous failure detection and gait recovery under stochastic, non-structured terrain conditions.
% +\revblue{The arbitration network autonomously transitions through this sequence.} The critical event occurs in Frames 5--6: rolling-mode wheel slip on the stochastic rough segment triggers a sustained stagnation signal, resolved by reactively redeploying the quadrupedal gait, which restores the support polygon and completes traversal. This closed-loop reactive switching was observed across all 15 trials, demonstrating reliable, autonomous failure detection and gait recovery under stochastic, non-structured terrain conditions.
% @@ -906 +873 @@
% -This subsection presents a rigorous quantitative analysis of the proposed bio-inspired controller evaluated on the physical hardware platform across four distinct stimulus-driven scenarios, complemented by global energetic and latency profiling under variable terrain conditions. Three independent hardware replicates were executed per scenario family, and the resulting temporal, neural, and kinematic metrics were compiled into publication-quality analytical plots and consolidated in Table~\ref{tab:consolidated_physical_metrics} to provide full statistical transparency.
% +\revblue{This subsection reports four stimulus scenarios on the physical platform (three hardware replicates each, consolidated in Table~\ref{tab:consolidated_physical_metrics}), followed by energy and latency profiling over variable terrain.}
% @@ -975 +942 @@
% -\textbf{Complex (see Figure~\ref{fig:complex_real_compound}):} Under chained, safety-critical transformations over an extended hardware timeline, the deterministic step response shows that, across successive back-to-back transformations, the ROS~2 Control architecture effectively manages electromechanical inertia and trajectory blending, protecting the joints from cumulative torque stress. The three-phase firing-variance profile maps the multi-stage decision process: a transient hyper-excitation peak during static obstacle avoidance, driven by the physical barrier's high spatial density flooding the camera sensor; a secondary activation wave capturing the emergency flight reflex to the aversive cue; and a low-energy steady-state as the robot tracks the appetitive target to mission completion. The chassis attitude telemetry mirrors this sequence, with asymmetric flanking maneuvers exciting the transversal axis and rapid backward accelerations peaking the longitudinal tilt, before both metrics drop and stabilize during the final convergence phase, indicating the wheel-driven configuration smooths parasitic oscillations during target approach.
% +\textbf{Complex (see Figure~\ref{fig:complex_real_compound}):} \revblue{Across successive back-to-back transformations, the deterministic step response shows the ROS~2 Control architecture managing electromechanical inertia and trajectory blending, protecting the joints from cumulative torque stress.} The three-phase firing-variance profile maps the multi-stage decision process: a transient hyper-excitation peak during static obstacle avoidance, driven by the physical barrier's high spatial density flooding the camera sensor; a secondary activation wave capturing the emergency flight reflex to the aversive cue; and a low-energy steady-state as the robot tracks the appetitive target to mission completion. The chassis attitude telemetry mirrors this sequence, with asymmetric flanking maneuvers exciting the transversal axis and rapid backward accelerations peaking the longitudinal tilt, before both metrics drop and stabilize during the final convergence phase, indicating the wheel-driven configuration smooths parasitic oscillations during target approach.
% @@ -998 +965 @@
% -\textbf{Obstacle (see Figure~\ref{fig:obstacle_real_compound}):} Using physical barriers in the laboratory environment, the instantaneous step response, consistent across all three replicates, confirms the ROS~2 Control architecture processes physical geometric blockages as immediate emergency reflexes, bypassing trajectory ramp filters to guarantee instant hot-swapping of joint matrices. The massive transient spike in firing variance---an order of magnitude above the color-cue scenarios---reflects the volumetric obstacle flooding the camera sensor with spatial and depth data, triggering a hyper-excitation wave that drives aggressive lateral flanking, collapsing back to baseline once the chassis clears the obstacle's spatial boundary. The rhythmic Roll expansion with coupled Pitch variation reflects the weight-shifting needed for a tight lateral bypass, flattening once safety is secured.
% +\textbf{Obstacle (see Figure~\ref{fig:obstacle_real_compound}):} \revblue{As in the Aversive scenario, the step response is immediate and consistent across all three replicates: geometric blockages are handled as emergency reflexes without trajectory ramp filtering.} The massive transient spike in firing variance---an order of magnitude above the color-cue scenarios---reflects the volumetric obstacle flooding the camera sensor with spatial and depth data, triggering a hyper-excitation wave that drives aggressive lateral flanking, collapsing back to baseline once the chassis clears the obstacle's spatial boundary. The rhythmic Roll expansion with coupled Pitch variation reflects the weight-shifting needed for a tight lateral bypass, flattening once safety is secured.
% @@ -1025 +992 @@
% -The energetic cost of physical adaptation was quantified from individually detected structural-command transitions \revblue{as described in Figure~\ref{fig:MorphologicalEnergy}'s caption}. The Quadrupedal $\rightarrow$ Differential Rolling transition required $5.06 \pm 2.93$~J ($N=16$), and the reverse transition (Differential Rolling $\rightarrow$ Quadrupedal) required $4.13 \pm 1.19$~J ($N=19$); the between-direction difference is not statistically significant \revblue{(see caption)}, so neither direction should be read as more energetically costly than the other. Greater trial-to-trial variability was observed in the Quadrupedal $\rightarrow$ Differential Rolling direction, though the present dataset does not isolate a specific mechanical cause for this dispersion. As the Quadrupedal$\leftrightarrow$Differential Rolling pair involves the largest morphological reconfiguration among the three modes (full limb deployment/retraction together with the leg-to-wheel locomotion change), its transition energy serves as an upper bound on the cost of any mode transition: in the worst case a transition consumes $\leq 5.06$~J, a small cost relative to continuous operation (e.g.\ a typical 38\,s Omnidirectional segment integrates to $\approx 216$~J). To directly address the reviewer's request for long-term average energy consumption across locomotion modes, Table~\ref{tab:mode_steady_state_energy} reports the mean steady-state power draw for each of the three modes, computed from pooled intra-segment power samples with the electromechanical coupling window excluded from each mode segment detected via the robot's structural-command signal.
% +The energetic cost of physical adaptation was quantified from individually detected structural-command transitions \revblue{as described in Figure~\ref{fig:MorphologicalEnergy}'s caption}. The Quadrupedal $\rightarrow$ Differential Rolling transition required $5.06 \pm 2.93$~J ($N=16$), and the reverse transition (Differential Rolling $\rightarrow$ Quadrupedal) required $4.13 \pm 1.19$~J ($N=19$); the between-direction difference is not statistically significant \revblue{(see caption)}. Greater trial-to-trial variability was observed in the Quadrupedal $\rightarrow$ Differential Rolling direction, though the present dataset does not isolate a specific mechanical cause for this dispersion. As the Quadrupedal$\leftrightarrow$Differential Rolling pair involves the largest morphological reconfiguration among the three modes (full limb deployment/retraction together with the leg-to-wheel locomotion change), its transition energy serves as an upper bound on the cost of any mode transition: in the worst case a transition consumes $\leq 5.06$~J, a small cost relative to continuous operation (e.g.\ a typical 38\,s Omnidirectional segment integrates to $\approx 216$~J). \revblue{Table~\ref{tab:mode_steady_state_energy} reports the long-term average steady-state power draw of each mode.}
% @@ -1065 +1032 @@
% -As an open, low-cost exploratory research baseline, the physical robot is constructed using low-torque servomotors and a lightweight structural layout, which limits its operational envelope compared to industrial-grade platforms. Locomotion gaits are currently executed via pre-configured open-loop control patterns; while this reduces control complexity, it weakens active stabilization capabilities on high-slip surfaces or excessively steep inclinations. This constraint is a direct consequence of the platform's low-cost mechanical design ---low-torque servomotors and a lightweight PLA structure--- rather than a limitation of the neural arbitration architecture itself. As a representative case, the gait selected during steep-slope negotiation (Figure~\ref{fig:RES_IMU2_real}) reflects the boundaries of the physical actuator envelope: the neural decision module consistently selects the locomotion mode associated with improved stability under the detected terrain conditions, yet the open-loop execution of that gait cannot compensate for the torque deficits inherent to the hardware. Experimental validation was constrained to structured laboratory environments with static stimuli and did not include unconstrained outdoor topographies, dynamic obstacles, or severe external disturbances. Therefore, the empirical results presented herein do not constitute validation of a production-ready system for rugged industrial deployment. Instead, they substantiate the mathematical, biological, and computational feasibility of a highly efficient basal-ganglia inspired arbitration network capable of managing multimodal mobility on resource-constrained embedded hardware.
% +As an open, low-cost exploratory research baseline, the physical robot is constructed using low-torque servomotors and a lightweight PLA structure, which limits its operational envelope compared to industrial-grade platforms. Locomotion gaits are currently executed via pre-configured open-loop control patterns; while this reduces control complexity, it weakens active stabilization capabilities on high-slip surfaces or excessively steep inclinations. \revblue{This constraint stems from the hardware, not from the neural arbitration architecture.} As a representative case, the gait selected during steep-slope negotiation (Figure~\ref{fig:RES_IMU2_real}) reflects the boundaries of the physical actuator envelope: the neural decision module consistently selects the locomotion mode associated with improved stability under the detected terrain conditions, yet the open-loop execution of that gait cannot compensate for the torque deficits inherent to the hardware. Experimental validation was constrained to structured laboratory environments with static stimuli and did not include unconstrained outdoor topographies, dynamic obstacles, or severe external disturbances. Therefore, the empirical results presented herein do not constitute validation of a production-ready system for rugged industrial deployment. Instead, they substantiate the mathematical, biological, and computational feasibility of a highly efficient basal-ganglia inspired arbitration network capable of managing multimodal mobility on resource-constrained embedded hardware.
% @@ -1068 +1035 @@
% -\revblue{Regarding the computing platform: the neural circuits, the MLP terrain classifier, and the image processing ran on an external laptop linked to the robot through a serial connection (Section~\ref{ssec:HardwareSpec}), because the onboard ESP32-S3 is dedicated to sensor acquisition and actuator control. Onboard feasibility was assessed only for the MLP classifier (Section~\ref{sssec:mlp_eval}); the basal-ganglia, locomotion, and gait-decision circuits and the vision pipeline were not benchmarked on the microcontroller, so a fully autonomous onboard implementation remains future work.}
% +\revblue{Regarding the computing platform: the neural circuits, the MLP terrain classifier, and the image processing ran on the external laptop of Section~\ref{ssec:HardwareSpec}, not on the robot's ESP32-S3. Onboard feasibility was assessed only for the MLP classifier (Section~\ref{sssec:mlp_eval}); the basal-ganglia, locomotion, and gait-decision circuits and the vision pipeline were not benchmarked on the microcontroller, so a fully autonomous onboard implementation remains future work.}
% @@ -1074,2 +1041 @@
% -of an industrial storage facility, included to contextualize
% -the platform's application domain.
% +of an industrial storage facility.
% --- a/sections/conclusions.tex
% +++ b/sections/conclusions.tex
% @@ -16 +16 @@
% -The findings provide a highly efficient baseline demonstrating that lightweight, bio-inspired neural networks ---arbitrating between locomotion modes and reactive behaviors driven by visual and inertial cues--- can manage multimodal gait transitions \revblue{on a real robotic platform, with the neural circuits executed on an external computer}, allowing it to navigate, stabilize, and make decisions robustly in environments with obstacles and adverse terrains. The ability to switch locomotion modes based on inertial stability metrics expands the robot's operational range across flat, rough, and inclined surfaces, while the integration of camera, LiDAR, and IMU into low-latency decision loops enabled dynamic goal prioritization, collision avoidance, and timely, collision-free goal approach, validated in both simulation and physical experiments. The perception-decision pipeline additionally retained task-level performance under combined IMU-drift and LiDAR-noise perturbation, and, conditionally on the scenario, camera illumination distortion (Section~\ref{ssec:NoiseRobustnessEval}), evaluated jointly rather than as isolated single-sensor degradations. When scaled onto robust physical platforms, this low-cost computational paradigm illustrates how decentralized, bio-inspired architectures could optimize energy efficiency and behavioral reliability, serving as a stepping stone toward bridging open-loop exploratory prototypes and the rigorous stabilization required by unconstrained, real-world deployment in power-plant monitoring or automated industrial operations. Although the physical prototype displays hardware-level boundaries ---including low-torque actuation and open-loop locomotion profiles--- these results substantiate that complex, multi-stimulus mobility decisions can be \revblue{taken for such resource-constrained platforms; onboard feasibility was assessed only for the MLP terrain classifier on the ESP32-S3 (Section~\ref{sssec:mlp_eval}), not for the full neural pipeline}, with the future-work priorities detailed below.
% +\revblue{The findings show that lightweight, bio-inspired neural networks ---arbitrating between locomotion modes and reactive behaviors driven by visual and inertial cues--- can manage multimodal gait transitions on a real robotic platform, with the neural circuits executed on an external computer.} The perception-decision pipeline additionally retained task-level performance under combined IMU-drift and LiDAR-noise perturbation, and, conditionally on the scenario, camera illumination distortion (Section~\ref{ssec:NoiseRobustnessEval}), evaluated jointly rather than as isolated single-sensor degradations. \revblue{Although the physical prototype displays hardware-level boundaries (Section~\ref{ssec:Limitations}),} these results substantiate that complex, multi-stimulus mobility decisions can be \revblue{taken for such resource-constrained platforms; onboard feasibility was assessed only for the MLP terrain classifier on the ESP32-S3 (Section~\ref{sssec:mlp_eval}), not for the full neural pipeline}.
